Alexander is a climate change denier and wants to prove to Joshua that the Earth is not getting warmer at all.
Alexander argues that it has snowed several times already in November, suggesting that the temperature is in fact dropping.
Joshua, who mostly stays indoors, hasn't noticed any snow. 
Therefore, Alexander has been running around taking pictures of particles in the air, 
hoping to provide photographic evidence to Joshua. 
Now he has finally got an image he is quite satisfied with, 
but he needs to make some changes before showing the picture.

Joshua has certain perceptions of how a snowflake should look. 
An image consists of $N \times N$ pixels, where $N$ is an odd number, 
and each pixel is either white or black. 
The snowflake is composed of the white pixels in the image. 
According to Joshua, the image represents a snowflake if the following criteria are met:


\begin{itemize}
\item The center pixel is white.
\item If you rotate the image 90 degrees clockwise, it looks the same.
\item If you mirror the image along the vertical centerline, it looks the same.
\item The set of white pixels is contiguous \emph{without diagonal steps}. 
  This means you can draw a curve from any white pixel to any other white pixel without the pen passing through a black pixel. 
  The pen can move up, down, left, and right, but not diagonally.
\item The set of black pixels is contiguous \emph{with diagonal steps}. 
This means you can draw a curve from any black pixel to any other black pixel without the pen passing through a white pixel. 
  The pen can move up, down, left, right, \emph{and} diagonally.
  \item There are no hooks in the snowflake. Let $\text{pixel}\text{up}$, $\text{pixel}\text{right}$, $\text{pixel}\text{down}$, and $\text{pixel}\text{left}$ be the adjacent pixels to a white pixel. A hook occurs if exactly two of these adjacent pixels are also white, and they are not on opposite sides of the pixel. Therefore, for example, if $\text{pixel}\text{up}$ and $\text{pixel}\text{right}$ are white while $\text{pixel}\text{down}$ and $\text{pixel}\text{left}$ are black, a hook is formed. However, if $\text{pixel}\text{up}$ and $\text{pixel}\text{down}$ are white, and $\text{pixel}\text{right}$ and $\text{pixel}\text{left}$ are black, a hook does not form.

\end{itemize}

Alexander has downloaded an image editing program where he can change the color of pixels, 
both from black to white and from white to black. 
Since he is not very skilled at image editing, Alexander expects it to take him one second for each pixel he needs to change. 
Alexander is a very busy man and wants to be as time-efficient as possible. 
He now wants to know how quickly he can finish if he changes as few pixels as possible.

\textbf{Input:}
The first line of the input contains an integer $N$ $(7 \le N \le 15)$, the number of types of buildings.

Following that are $N$ lines, where the $i$-th line contains the $N$ numbers $p_{i,1}, p_{i,2}, ..., p_{i,N}$. If the pixel at position $i, j$ in the image is white, then $p_{i,j} = 1$; otherwise, $p_{i,j} = 0$.

\textbf{Output:}
Output an integer: the number of seconds it takes for Alexander to make changes to his image so that Joshua agrees that it represents a snowflake.

\textbf{Scoring:}
Your solution will be tested on a set of test case groups. To earn points for a group, you must pass all test cases in that group.

\begin{tabular}{| l | l | p{12cm} |}
\hline
\textbf{Group} & \textbf{Points} & \textbf{Limits} \ \hline
$1$ & $20$ & $N=7$. \ \hline
$3$ & $20$ & $N=11$ \ \hline
$4$ & $60$ & $N=15$ \ \hline
\end{tabular}